\answer
\begin{enumerate}

%% 9章 問1（9.2節）
\item[9.2節 問1]
$v_R$は$+100$\,Vと$-100$\,Vを行き来する方形波で振幅$100$\,V，
抵抗負荷の電流は$i_R = v_R/R$より振幅$100/10 = 10$\,Aの同相の方形波である。
負荷を誘導性(RL)に替えると，電流は方形波電圧を積分した形になり，
各半周期で指数関数状に増減する三角波状の波形になる。
またインダクタが電流の変化を妨げるため，電流の位相は電圧より遅れる。

%% 9章 問2（9.2節）
\item[9.2節 問2]
$v_R$を$+E$から$-E$へ切り替えた直後は，$\mathrm{S}_2$と$\mathrm{S}_3$をオンにしたい状態だが，
電流はまだ$+$の向き（左端から負荷を通って右端へ）に流れ続けている。
この電流はデッドタイム中，右レグ下アームの還流ダイオード（$\mathrm{S}_4$に並列）を通って下レールへ入り，
左レグ上アームの還流ダイオード（$\mathrm{S}_1$に並列）を通って戻る。
すなわち\textbf{$\mathrm{S}_4$と$\mathrm{S}_1$に並列な2つの還流ダイオード}を電流が通る。
これらは電源へエネルギーを戻す向きに導通している。

%% 9章 問3（9.3節）
\item[9.3節 問3]
式\ref{eq9.8}より高調波の振幅は基本波の$1/n$だから，
3次は$1/3 \approx 33$\,\%，5次は$1/5 = 20$\,\%である。
基本波の実効値は式\ref{eq9.11}より$V_1 = 4E/(\pi\sqrt{2}) = 2\sqrt{2}E/\pi \approx 0.90E$。
方形波全体の実効値は（振幅$E$で常に$\pm E$だから）$E$なので，
基本波は全体の約$90$\,\%を占める。

%% 9章 問4（9.4節）
\item[9.4節 問4]
$\cos^2(\omega\mathit{\Delta}t/2) = 1/2$より$\cos(\omega\mathit{\Delta}t/2) = 1/\sqrt{2}$，
すなわち$\omega\mathit{\Delta}t/2 = \pi/4$である。
$\omega = 2\pi/T$を代入すると
\begin{equation*}
\frac{2\pi}{T}\cdot\frac{\mathit{\Delta}t}{2} = \frac{\pi}{4}
\quad\Rightarrow\quad
\mathit{\Delta}t = \frac{T}{4}
\end{equation*}
半周期$T/2$の半分だけずらせば，基本波は方形波の半分になる。

%% 9章 問5（9.5節）
\item[9.5節 問5]
基本波の実効値が$100$\,Vだから振幅は$\hat{V}_1 = 100\sqrt{2} \approx 141$\,V。
ハーフブリッジは式\ref{eq9.16}より$\hat{V}_1 = mE/2$なので
\begin{equation*}
m = \frac{2\hat{V}_1}{E} = \frac{2 \times 141}{300} \approx 0.94
\end{equation*}
$m = 0.94 < 1$だから線形制御の範囲に収まっている。

%% 9章 問6（9.6節）
\item[9.6節 問6]
3レベル（中性点クランプ）方式では，オフのスイッチにかかる電圧は
電源電圧$E$の\textbf{半分}（$E/2$）である。
2レベルではオフのスイッチに$E$がまるごとかかるのに対し，
3レベルでは中点にクランプされて$E/2$で済む。
このため，同じ電源電圧でも耐圧の低い（＝安価でオン抵抗の小さい）素子が使え，
高電圧のインバータでは1素子では耐えられない電圧を分担して扱えるようになる。

%% 9章 問7（9.7節）
\item[9.7節 問7]
\begin{equation*}
I = \frac{Q_g}{t} = \frac{65\times10^{-9}}{100\times10^{-9}} = 0.65\,\mathrm{A}
\end{equation*}
単位は
\begin{equation*}
\frac{\mathrm{C}}{\mathrm{s}} = \frac{\mathrm{A\cdot s}}{\mathrm{s}} = \mathrm{A}
\end{equation*}
となり，確かに電流の単位になる（$1\,\mathrm{A} = 1\,\mathrm{C/s}$）。
$0.65$\,Aはマイコンの出力ピンには荷が重く，専用のゲートドライバが必要であることがわかる。

\end{enumerate}
